% ============================================================
% 鉴历史·文采说（杂志页码 30–33）
% ============================================================
\colchapter{鉴历史·文采说}

\section*{嘲讽过后，鲁迅先生相信怎样的英雄？}
\auth{文/张佳玮}

鲁迅先生的笔触向来是锋利且冷峻的，在蒙昧、落后的旧中国，那就仿佛一柄利剑，刺破了压迫之下阻碍百姓思想启迪的层层乌云。只不过，在生命最后的时光里，先生同样塑造了最终拥有喜剧结局的大禹，他是一个脚踏实地服务人民并真正挺起脊梁的英雄，也是先生对于黑暗终将过去，光明必将到来的深情期盼。

鲁迅先生的小说里，很少写英雄有好结局。他笔下惯多苦人如祥林嫂、闰土；乡民如阿Q、九斤老太……各有各的苦。而算是大英雄且有喜剧结尾的，似乎只有一篇——《理水》，讲的是大禹治水。

《理水》的世界里，天下正遭大水灾。然而，百姓混沌，学者萎靡。学者嗔怪乡下人没文化，绅士出来责怪百姓，说大洪水都是百姓弄出来的：“水还没来的时候，他们懒着不肯填，洪水来了的时候，他们又懒着不肯戽……”而专员们来考察，玩了一圈，见了个百姓代表，也没反映啥。回京之后，考察员们大排筵宴，接风洗尘，列鼎吃肉……

终于，大禹出现了：没穿袜子，满脚底都是栗子一般的老茧。他听了专家们不把百姓当人的胡言乱语，说了很实在的话，“我经过查考，知道先前的方法：‘湮’，确是错误了。以后应该用‘导’！不知道诸位的意见怎么样？”当然，立刻有官员反对：“这是蚩尤的法子！”有个白发老官还跟大禹强调：“湮”是世界上已有定评的好法子，别的现代方法，都有问题，不能用。另外，有人强调大禹父亲以前犯过的过错之类。

大禹微微一笑：“我知道的。有人说我的爸爸变了黄熊，也有人说他变了三足鳖，也有人说我在求名，图利。说就是了。我要说的是我查了山泽的情形，征了百姓的意见，已经看透实情，打定主意，无论如何，非‘导’不可！这些同事，也都和我同意的。”他所谓的同事，是一群黑瘦的乞丐似的人：不动，不言，不笑，像铁铸的一样。他们没有专家大员学者那么能忽悠，只是在实地做事，做得一身黑瘦、两脚老茧……最后，大禹治水成功，去见了舜。

我们都知道鲁迅先生爱嘲弄各色人等，古来圣贤在《故事新编》里多少都落了窘境。相比起来，大禹这个形象，是先生心里理想的英雄角色了，最后也的确大获成功，走向大团圆结局。鲁迅笔下的大禹，说来无非是：实地考察做实事，为做出正确抉择，不惜力排众议；懒得说好听的，每天只琢磨着受苦的百姓。他对舜并不歌功颂德，只平静地说：车船橇轿，到山砍树，放水入川；吃饭吃肉，搬家安居。而他给舜的建议也很直白：“要小心、安静。对天有良心。”这样的人物，即便是黑瘦的、乞丐似的，脚下起茧子，也是真正的英雄。

鲁迅先生写《理水》时，距离逝世，已经只有十一个月了。他看过所有的黑暗，经历过所有的战斗，完成了人生绝大多数的嘲讽。许多对手说他悲观，许多对手说他刻薄，也有许多对手说他越到后期越“毒舌”。的确，那个时代暗到了一个程度，但到最后，大概鲁迅先生还是相信：这样实实在在做实事、不整花里胡哨虚招子、把百姓切实生活放心里的“小心、安静、有良心”的人，才是华夏真正的脊梁。

\begin{quoteBox}{文采沙龙}
从《呐喊》到《彷徨》，鲁迅先生一生都在思考国民性与民族出路的问题。而他笔下的“理水”，隐喻的正是那个灾难频仍的中国。在那种境况下，先生依然选择相信大禹式脚踏实地的实干家，相信“小心、安静、有良心”的朴素道义——这是何等深沉而坚定的信念。
\end{quoteBox}

\clearpage
\section*{“世界小姐”光环散去：美丽，不需要模板}
\auth{文/叶橙子}

时间倒回十多年前，那时的人们或许很难想象今日“世界小姐”评选的尴尬局面：这档世界选美节目，曾经的全球电视观众多达20亿；如今在国内，却陷入了中国区候选人学历造假的丑闻。更荒唐的是，网友们扒到最后，发现连主办方都出现了真假两个。围观群众则更是疑惑——“这比赛原来还在办吗？”

绝大多数人对“世界小姐”评选的记忆，都停留在舞台上各国美女身着泳衣或礼服款款走动，微微挥手，披好“MISS WORLD”的绶带，再戴上一顶皇冠……“世界小姐”们温婉自信、在电视那头遥不可及。但如今，这项大赛的光环不再稀有——不仅在网红视频里随时可见，这一次还成了互联网的围观笑料。

曾经的“世界小姐”评选想办法抹平着地域间的审美差异。其所依靠的手段，是巧妙地提炼出一种当年的世界完美女性的公约数。在一众选美比赛当中，“世界小姐”是考核程序上看起来最为“德智体美劳”全面发展的一个——除外貌身材的考核，还考查模特的才艺、体育、心灵与智慧美等。说到底，“世界小姐”在单项上未必做到了最好，却面面俱到，呈现出一种平均的优秀。而人们则一同造一场乐见其成的幻梦，一同沉浸其中：选秀佳丽们在舞台上展现各国民俗特色，构建出地球村的样貌，宣传着世界和平与关爱孩子的口号。

一切都很美好，但这场幻梦的构建是有条件的。当外界改变，“世界小姐”的光环就再也遮掩不住选美活动本身的缺陷。就像“世界小姐”的活动宗旨里，有一项是“促进世界和平”，于是舞台上，佳丽们被询问如何看待恐怖主义、如何阻止全球变暖、如何帮贫困地区的孩子改善生活等等庞大的问题。但在短暂的考核中，她们其实只能模糊地回答“愿世界充满爱与和平”。由此，这一环节更像是某种理想主义的象征，而随着世界布局变化、情绪走向撕裂，舞台上的象征表演开始显得不合时宜。

一场选美比赛，要有地办、有人信、有观众看才行。但在近十几年内，“世界小姐”慢慢失去了三个领域的大部分优势。一来网络媒体逐渐占领生活之后，电视节目的影响力本就在世界范围内大大降低，让城市在“世界小姐”里露脸打广告，已然成为落伍的手段。再者，“世界小姐”等国际选美比赛爆火之后，相似比赛纷纷冒出，外界对“世界小姐”商业价值的判断愈发走低……

曾经，“世界小姐”的野心很大，试图用一种美丽去代表人群，也试图用一个人代表国家，来表演世界大同的温和场景。不过，现在看来，如今的人们已不再需要这些光鲜的表演，不需要一种被定义的美丽。而美，绝不止有一种解释，也绝不再是谁的唯一价值。

\begin{quoteBox}{文采沙龙}
“世界小姐”的选拔标准本质是四平八稳的优雅美人。优雅美人很难被推翻、被讨厌，但在当今滚滚向前的审美浪潮之中，这样的选拔会显得无趣而平淡。冠军们的挥手与回眸仿佛是一个模子刻出的。世界范围内的女性都需要更多关于“美”的展现，需要更多魅力与力量的方式，独独不需要一个被循规蹈矩的名号。
\end{quoteBox}

\clearpage
\section*{用对“的地得”：跨越千年，致敬中国文字}
\auth{文/王申中}

新学年，对于师生们来说就像新年之始。一位小学语文教师这样许下自己的“新年愿望”：祈祷期末考试不要出现“的地得”，不愿再考，“跟学生携手斗了6年，终究还是斗不过它们。”“的地得”，不仅是小朋友学不好，好些成年人也学不会。中文互联网苦乱用“的地得”久矣。

平心而论，正确使用“的地得”确实不是一项自古至今的历史传统。为了分清“的地得”，我们花了1000年的时间。使用频率最高的“的”，最开始其实是个形容词，指明亮、光明的样子。直到唐末，它才被用作助词，不过那时被写作“底”。宋元时期，民间小说兴起，大量口语进入文本，“底”才开始被“的”取代。而五四运动以后，汉语受欧化的影响，形成了“的、地、底”三足鼎立的局面，“的”“底”用在定语后面，状语后面用“地”。不过，这种局面到20世纪40年代被打破了，“底”逐渐被人们冷落，成了“的”“地”平分秋色。

那么，“得”去哪儿了？它一直在独自美丽。在甲骨文里就已经出场的“得”，许多年来一直被当作“获得、得到”的意思来使用。唐朝时，“得”就有了宾语用途，逐渐接近现代汉语意义上的助词“得”。

所以，词义起源本就不相同的“的地得”，怎么就成了现代人混用、乱用的重灾区？公众号“汉语堂”推测，虽然“的地得”出现超过千年，但一直没有官方定论。直到1956年，《暂拟汉语教学语法系统》才明确且权威地规定：“它们是定语、状语、补语的标志。”不过，1984年公布的《中学教学语法系统提要（试用）》又提出，在中小学教学中不再对“的地得”的用法做硬性区分。

用错“的地得”，还可以怪谁？怪上网冲浪。互联网时代，线上交流成了当代人最主流的社交模式，也诞生了“的地得”混用，标点符号乱用、不用，一句话分成几句话发等等特点。传播语境的碎片化打破了传统的话语交流模式，讲究即时即刻沟通的线上交流更追求短平快，错字或别字成为一种被允许的存在。特别是使用拼音输入法时，我们思考运用的是声音符号，而非汉字符号，导致了网络语言更接近口语而非书面语，相应地也降低了对用户书写字形的能力要求。

北宋僧人文莹编写的野史《湘山野录》中记录了这样一个故事。宋太祖赵匡胤经过朱雀门，看到城门上匾额写的是“朱雀之门”，就问随行的大臣赵普，“何不只书朱雀门，须著之字安用？”赵普答道：“语助。”意思是说这就是个助词。赵匡胤于是大笑说：“之乎者也，助得甚事？”笑话文人咬文嚼字，虚词能顶什么用啊。这也是成语“之乎者也”的来源。只不过，故事最后，赵匡胤还是不自觉地用了正确的“得”（助得甚事）。或许，这就是当今人们用好“的地得”的最大意义：对文字的尊重。

\begin{quoteBox}{文采沙龙}
当前，对于“的地得”的误用，其实已无碍人们的正常交流，但正是在这样一种心照不宣的使用场景之下，我们更应注意维护中国文字表情达意的严肃性。毕竟，走过1000年的历史，一个汉字、一种表达才得到科学系统的沉淀，它们绝不该在社交狂欢的时代车轮下沦为可有可无、一笔带过的看客。
\end{quoteBox}

\clearpage
\section*{“最后的女酋长”：林深时见鹿}

2022年8月20日，有着“中国最后一位女酋长”之称的玛利亚·索去世，享年101岁。如果你读过迟子建的小说《额尔古纳河右岸》，或是看过《歼达罕》《敖鲁古雅养鹿人》等纪录片，那么对于以养驯鹿为生的鄂温克族应该不会陌生。而玛利亚·索就曾带领这支中国最后的狩猎部族，生活在内蒙古大兴安岭深处。

每当有人问起玛利亚·索今年多大时，她总会回答“80岁”，年年如此。其实，她也不清楚自己的确切生日，在大兴安岭深处生活了一辈子，她没有现代人的时间概念，没用过日历、手表和电子产品，但她依然保持着狩猎部族的敏锐，靠自然的种种痕迹捕捉时间的答案。

玛利亚·索出生在激流河边，80岁前都没有离开过。她和其他鄂温克人一样，从小与驯鹿相伴——妈妈把她放在树皮做的筐子里，在驯鹿背上跟着大人一起搬家。几千年来，鄂温克人都遵循着这种古老的游牧规则，当驯鹿吃完这里的苔藓，他们就会收起撮罗子（帐篷），找到下一处安居之所。

鄂温克人与驯鹿是共生的关系。驯鹿温顺，能驮着自己的所有家当迁徙，也能载人；鹿奶好喝，也能揉进面团做成易保存的“鄂温克列巴”，是上山打猎必备。玛利亚·索曾回忆说，过去的人从来不杀驯鹿，也不吃，驯鹿死了都是风葬。然而，到了后来，有的族人开始吃鹿肉，玛利亚·索想都不敢想，“自家养的东西，怎么能舍得吃呢？”

丈夫去世后，玛利亚·索便独自扛起了家族的重担，族人都非常敬重她。外界称她为酋长，“中国最后一位女酋长”。但事实上，这些部落的酋长制度在1761年就被沙俄废除了，玛利亚·索更准确的身份应该是“族长”。2003年，玛利亚·索的部落迎来了历史的转折——在当地政府主导下，他们进行了一次大迁徙，从深林的猎民点迁出，来到百公里外的根河市郊。那时，大部分族人都住进了现代化小楼，玛利亚·索却无法接受这样的新生活，选择留在了阿龙山深处的猎民点。只是，鄂温克人打猎、游牧的生活终究是一去不返了。而随着年纪越来越大，玛利亚·索的身体状况不如以往，一只眼睛失明，耳背严重。女儿接她下山住的时间越来越久，尤其是在零下30多摄氏度的冬天。

就这样，在人生最后二十年里，玛利亚·索的精神世界被时间割成两半，一半留在了山上，继续与驯鹿做伴，另一半则被现代化发展的车轮撵着，不得不往前。她时常会想起年轻的时候，鹿群在头鹿带领下，钻进山林找苔藓，一去就是十几公里。鄂温克人进山找鹿，“喂——来来来”地一叫，然后站在原地，闭起眼睛，驯鹿脖子上的铃铛就会在远处“叮叮当当”地响起。如今，玛利亚·索走了，山林里驯鹿的铃铛声也随之消散。

\begin{quoteBox}{文采沙龙}
玛利亚·索老人之所以总回答自己是“80岁”，或许不仅因为记忆的模糊，也是由于她把自己最幸福的时光都留在了那个年纪。那时，她仍然可以在山林中和驯鹿一起自由嬉戏，没有面临城市生活的侵扰，而自此之后，她只能用力地维持一份旧时光的惯性，并不断与时代的变化拉扯，目送驯鹿们的背影渐行渐远。
\end{quoteBox}
